\documentclass[11pt,a4paper]{article}
\usepackage[T1]{fontenc}
\usepackage{lmodern}
\usepackage{amsmath,amssymb,amsthm}
\usepackage[margin=25mm]{geometry}
\usepackage[hidelinks]{hyperref}

\DeclareMathOperator{\Tr}{Tr}
\newcommand{\trv}{\underline{\Tr}}
\newcommand{\R}{\mathbb{R}}
\newcommand{\C}{\mathbb{C}}
\newcommand{\ind}{\mathbf{1}}
\newtheorem*{theorem}{Theorem}
\setlength{\parindent}{0pt}
\setlength{\parskip}{0.4em}

\begin{document}
\begin{center}
  {\Large\bfseries Graphene conductivity}
\end{center}
The goal is to establish the residual conductivity of graphene, within
linear response, in a non-interacting model.

\textbf{Setting.}
Spinless electrons, atomic units.
Let $\mathcal R\subset\R^2$ be a lattice, $\Omega$ a unit cell, and
$\mathcal B$ a unit cell of
\[
  \mathcal R^*=\{K\in\R^2:K\cdot R\in2\pi\mathbb Z
  \text{ for every }R\in\mathcal R\}.
\]
For real $V\in H^1_{\mathrm{per}}(\Omega)$, set
\[
  H=-\tfrac12\Delta+V\quad\text{on }L^2(\R^2;\C),
  \qquad
  H_k=\tfrac12(-i\nabla+k)^2+V
  \quad\text{on }L^2_{\mathrm{per}}(\Omega;\C).
\]
$H_k$ is self-adjoint with domain
$H^2_{\mathrm{per}}(\Omega;\C)$. Write its eigenvalues, with multiplicities, as
$\lambda_1(k)\leq\lambda_2(k)\leq\cdots$, and set
$\lambda_0(k)=-\infty$.

\textbf{Band assumptions.}
Fix $\mu_{\mathrm F}\in\R$, and assume that the system is semi-metallic: there are
distinct points
$k_1,\ldots,k_M$ modulo $\mathcal R^*$ and a number of bands $N$ such that
\[
  \lambda_N(k)<\mu_{\mathrm F}<\lambda_{N+1}(k)
  \qquad(k\notin\{k_1,\ldots,k_M\}\bmod\mathcal R^*),
\]
and, for each $i=1,\ldots,M$,
\[
  \lambda_{N-1}(k_i)<\lambda_N(k_i)=\mu_{\mathrm F}
  =\lambda_{N+1}(k_i)<\lambda_{N+2}(k_i).
\]
There are $v_i>0$ such that, as $q\to0$,
\[
  \lambda_N(k_i+q)=\mu_{\mathrm F}-v_i|q|+O(|q|^2),
  \qquad
  \lambda_{N+1}(k_i+q)=\mu_{\mathrm F}+v_i|q|+O(|q|^2).
\]

\textbf{Current.}
Fix $e_{\beta},e_{\alpha}\in\R^2$, with $e_{\beta}$ the field direction
and $e_{\alpha}$ the current direction.
For all $\varepsilon\geq0$ and $t\geq0$, define
\begin{gather*}
  \gamma_0=\ind_{(-\infty,\mu_{\mathrm F}]}(H),
  \qquad H_\varepsilon=H-\varepsilon x\cdot e_{\beta},
  \qquad J_{\alpha}=(-i\nabla)\cdot e_{\alpha},\\
  \gamma_\varepsilon(t)
  =e^{-itH_\varepsilon}\gamma_0e^{itH_\varepsilon},
  \qquad
  j_\varepsilon(t)
  =\trv\bigl(J_{\alpha}\gamma_\varepsilon(t)\bigr),
\end{gather*}
where $H_\varepsilon$ denotes the self-adjoint closure of its
restriction to $C_c^\infty(\R^2;\C)$, and
\[
  \trv(A)=\frac1{|\Omega|}\Tr(\ind_\Omega A\ind_\Omega)
\]
for locally trace-class periodic $A$.

\begin{theorem}
Under the above assumptions, the iterated limit exists and
\[
\lim_{T\to+\infty}\frac1T\int_0^T
  \left(\lim_{\varepsilon\to0^+}
  \frac{j_\varepsilon(t)}{\varepsilon}\right)\,dt
  =\frac{M}{16}\,e_{\alpha}\cdot e_{\beta}.
\]
\end{theorem}
\end{document}
